% verso del frontespizio: colophon
% Adattato dall'originale del template, che era impostato per una tesi di
% dottorato (imprint con copyright e "Published by").
\thispagestyle{empty}
\hphantom{.}
\vfill

\section*{Colophon}
\textit{\TitleText}\\
Tesi di laurea triennale in Informatica, Universit\`a degli Studi di Camerino,
anno accademico 2025--2026.\\[0.5\baselineskip]
%
Il documento \`e composto in \LaTeX{} con la classe \texttt{unicam-diss} di Matteo
Belenchia, basata su \texttt{tufte-book}\footnote{\url{https://github.com/Tufte-LaTeX/tufte-latex}}
e sullo stile \texttt{classicthesis}\footnote{\url{https://bitbucket.org/amiede/classicthesis/}}
di Andr\'e Miede.\\[0.5\baselineskip]
%
La bibliografia \`e elaborata con Biblatex e Biber. Il carattere con grazie \`e
XCharter di Michael Sharpe, estensione del Bitstream Charter; il carattere senza
grazie \`e \TeX\ Gyre Heros; il monospaziato \`e Bitstream Vera Mono di Jim Lyles.
I titoli usano il Concrete Modern di Knuth.\\[0.5\baselineskip]
%
I dati sperimentali sono stati acquisiti con firmware e script scritti per questo
lavoro (Appendice~\ref{app:codice}); i file CSV originali delle sessioni di misura
sono conservati insieme ai sorgenti del documento.
